\documentclass[UTF8,a4paper,11pt,fontset=fandol]{ctexart}
\usepackage[margin=24mm]{geometry}
\usepackage{amsmath,amssymb,mathtools,stmaryrd}
\usepackage{booktabs,tabularx,array}
\usepackage{xcolor,listings,xurl,fancyhdr}
\usepackage[unicode,colorlinks=true,linkcolor=blue!45!black,citecolor=blue!45!black,urlcolor=blue!45!black]{hyperref}
\pagestyle{fancy}
\fancyhf{}
\fancyhead[L]{\small Pipeline Core Language}
\fancyhead[R]{\small 语法与操作语义}
\fancyfoot[C]{\thepage}
\setlength{\headheight}{14pt}
\setlength{\parskip}{0.35em}
\setlength{\emergencystretch}{3em}
\renewcommand{\arraystretch}{1.18}
\allowdisplaybreaks[2]
\newcolumntype{Y}{>{\raggedright\arraybackslash}X}
\newcommand{\code}{\nolinkurl}
\newcommand{\sem}[1]{\llbracket #1\rrbracket}
\newcommand{\kw}[1]{\mathsf{#1}}
\newcommand{\dom}{\operatorname{dom}}
\newcommand{\kind}{\operatorname{kind}}
\newcommand{\step}[1]{\xrightarrow{#1}}
\newcommand{\Read}{\operatorname{Reading}}
\newcommand{\Write}{\operatorname{Writing}}
\newcommand{\Compute}{\operatorname{Computing}}
\newcommand{\Done}{\operatorname{Done}}
\lstset{basicstyle=\small\ttfamily,columns=fullflexible,keepspaces=true,
  breaklines=true,showstringspaces=false,frame=single,
  rulecolor=\color{black!20},backgroundcolor=\color{black!2}}
\hypersetup{pdftitle={Pipeline Core Language：执行类别与并发操作语义},pdfauthor={}}
\title{Pipeline Core Language\\[0.4em]\large 执行类别与并发操作语义}
\author{}
\date{}

\begin{document}
\maketitle
\vspace{-2em}

\section{静态对象与执行类别}

令 $\mathsf{Pipe}$、$\mathsf{Addr}$、$\mathsf{Lock}$ 分别表示 pipeline、
存储单元地址和锁的标识集合。一个程序使用有限集合
$\mathcal P\subseteq\mathsf{Pipe}$。pipeline 标识包含所属设备、core 和执行
队列的身份。每个 pipeline 具有独立的顺序命令、局部变量和当前操作状态。
不同 pipeline 的步骤可以交错。

每个地址标识一个原子访问单元，有限地址集合表示访问区域。
共享内存 $\mu:\mathsf{Addr}\rightharpoonup\mathsf{Val}$ 保存地址和值。
源程序中具有别名的地址在翻译后必须保持相等。

执行类别由静态函数指定：
\[
\mathsf{ExecKind}=\{\mathsf{Cube},\mathsf{Vector},\mathsf{Scalar},\mathsf{MTE}\},
\qquad \kind:\mathsf{Pipe}\to\mathsf{ExecKind}.
\]
例如 $\kind(p_c)=\mathsf{Cube}$、$\kind(p_v)=\mathsf{Vector}$。
访问权限和函数支持关系由静态环境 $\Gamma$ 给出：
\[
\begin{aligned}
\operatorname{Supports}(p,f,\bar\tau)&\in\mathsf{Bool},&
\operatorname{CanRead}(p,a)&\in\mathsf{Bool},\\
\operatorname{CanWrite}(p,a)&\in\mathsf{Bool},&
\operatorname{TypeOK}(a,v)&\in\mathsf{Bool}.
\end{aligned}
\]
函数签名中的类型可以携带向量和矩阵形状。计算函数 $f$ 是定义在指定输入
类型上的确定性纯偏函数；其数值精度和舍入规则由具体实例给出。
计算函数不访问共享内存，也不操作事件和锁。

事件标识为 $\varepsilon=(s,p,q,k)$，分别表示同步域、源 pipeline、
目标 pipeline 和事件编号。$\Gamma$ 声明事件及每个同步域的参与者；
事件的两个端点必须属于相应同步域。每把锁 $\ell$ 声明允许参与的
pipeline 集合 $\operatorname{Participants}(\ell)\subseteq\mathcal P$。

\section{语法}

令 $x$ 为局部变量，$v$ 为值，$g$ 为纯表达式函数，$k$ 为执行类别。
值包括标量、布尔值、地址、有限地址集合、地址到值的有限映射、向量、
矩阵和 $\kw{unit}$。语法为：
\[
\begin{aligned}
e &::=v\mid x\mid g(e_1,\ldots,e_n),\\
o &::=\kw{read}(e_R)
  \mid\kw{write}(e_R,e_V)
  \mid\kw{computation}_{k,f}(e_1,\ldots,e_n),\\
a &::=x:=e\mid x:=o\mid\kw{setflag}(\varepsilon)
  \mid\kw{waitflag}(\varepsilon)
  \mid\kw{lock}(\ell)\mid\kw{unlock}(\ell),\\
C &::=\kw{skip}\mid a\mid C;C
  \mid\kw{if}\ e\ \kw{then}\ C\ \kw{else}\ C
  \mid\kw{while}\ e\ \kw{do}\ C,\\
P &::=\kw{parallel}\{p_1:C_1,\ldots,p_n:C_n\}.
\end{aligned}
\]
程序中各个 pipeline 标识互不相同，且
$\mathcal P=\{p_1,\ldots,p_n\}$。$\sem e_\rho$ 在局部环境 $\rho$ 中求值。
$e_R$ 求值得到有限地址集合 $R$；$e_V$ 求值得到定义域为 $R$ 的有限映射
$V$。读取返回地址值映射，写入返回 $\kw{unit}$。
例如 \code{unit := write(R,V)} 中的 \code{unit} 是接收返回值的局部变量。

同一 pipeline 的命令按程序顺序执行：操作完成、结果写入局部变量之后，
该 pipeline 才能执行后续命令。一个操作可以经历多个语义步骤，
这些步骤可以与其他 pipeline 的步骤交错。

\section{状态与操作初始化}

运行状态和 pipeline 局部状态分别为
\[
\Sigma=\langle\mu,F,L,T\rangle,\qquad
T(p)=\langle C_p,\rho_p,d_p\rangle.
\]
$F$ 将已声明事件映射到二值标志，$L$ 将已声明锁映射到空闲标记 $\bot$
或持有者 pipeline。$C_p$ 是剩余命令，$\rho_p$ 是局部环境，
$d_p$ 是空闲标记 $\bot$ 或当前操作的进度：
\[
d::=\Read(U,B)\mid\Write(U,V)\mid\Compute(k,f,\bar v)\mid\Done(v).
\]
$U$ 表示尚未访问的地址集合；$B$ 保存已经读取的值；$V$ 保存写入值。
执行计算期间保存执行类别、函数和已求值的参数。
$\Done(v)$ 表示内部执行已经完成，结果尚待写入局部环境。

在当前局部环境中初始化操作：
\[
\begin{aligned}
\operatorname{init}_\Gamma(\rho,\kw{read}(e_R))
 &=\Read(R,\varnothing)
 &&\text{当 }\sem{e_R}_{\rho}=R,\\
\operatorname{init}_\Gamma(\rho,\kw{write}(e_R,e_V))
 &=\Write(R,V)
 &&\text{当 }\sem{e_R}_{\rho}=R,\ \sem{e_V}_{\rho}=V,\\
\operatorname{init}_\Gamma(\rho,\kw{computation}_{k,f}(\bar e))
 &=\Compute(k,f,\bar v)
 &&\text{当 }\sem{\bar e}_{\rho}=\bar v.
\end{aligned}
\]
初始化后的状态还必须满足 $\operatorname{Valid}_\Gamma(p,\mu,d)$：
\[
\begin{aligned}
\operatorname{Valid}(p,\mu,\Read(R,\varnothing))
 &\iff R\subseteq\dom(\mu)\land
   \forall a\in R.\ \operatorname{CanRead}(p,a),\\
\operatorname{Valid}(p,\mu,\Write(R,V))
 &\iff R\subseteq\dom(\mu)\land\dom(V)=R\land\\
 &\quad\forall a\in R.\ \operatorname{CanWrite}(p,a)\land
   \operatorname{TypeOK}(a,V(a)),\\
\operatorname{Valid}(p,\mu,\Compute(k,f,\bar v))
 &\iff \kind(p)=k\land
   \operatorname{Supports}(p,f,\operatorname{types}(\bar v))\land
   f(\bar v)\text{ 有定义}.
\end{aligned}
\]
表达式求值失败、区域或数值类型错误、权限不足、函数未定义及计算类别
不匹配均产生错误。

初始状态中，$T(p)=(C_p,\rho_p,\bot)$，每个事件标志为 $0$，每把锁空闲。
初始内存符合 $\operatorname{TypeOK}$，初始局部环境只包含已声明的 pipeline；
$\Gamma$ 的事件、锁和计算函数声明满足上述条件。

\section{带标签的小步操作语义}

下列规则使用执行上下文
$E::=[\,]\mid E;C$。写作 $T[p\mapsto(C,\rho,d)]$ 的更新只改变 $p$
的局部状态；未列出的全局分量保持原值。标签中的
$\operatorname{start}$ 和 $\operatorname{complete}$ 记录一个操作的内部
执行步骤，不属于程序语法。

\subsection{顺序控制与操作}

局部赋值、顺序组合、条件和循环按以下规则执行：
\[
\frac{T(p)=(E[x:=e],\rho,\bot)\quad\sem e_\rho=v}
 {\Sigma\step{\operatorname{assign}(p,x,v)}
  \Sigma[T(p)\mapsto(E[\kw{skip}],\rho[x\mapsto v],\bot)]}
\]
\[
\frac{T(p)=(E[\kw{skip};C],\rho,\bot)}
 {\Sigma\step{\operatorname{control}(p)}
  \Sigma[T(p)\mapsto(E[C],\rho,\bot)]}
\]
\[
\frac{T(p)=(E[\kw{if}\ e\ \kw{then}\ C_1\ \kw{else}\ C_2],\rho,\bot)
 \quad\sem e_\rho=b}
 {\Sigma\step{\operatorname{control}(p)}
  \Sigma[T(p)\mapsto(E[C_b],\rho,\bot)]}
\]
其中 $C_{\kw{true}}=C_1$，$C_{\kw{false}}=C_2$。循环展开为
$\kw{if}\ e\ \kw{then}\ (C;\kw{while}\ e\ \kw{do}\ C)\
\kw{else}\ \kw{skip}$，并产生一个 $\operatorname{control}(p)$ 标签。

操作首先求值参数并检查合法性：
\[
\frac{T(p)=(E[x:=o],\rho,\bot)\quad
 \operatorname{init}_\Gamma(\rho,o)=d\quad
 \operatorname{Valid}_\Gamma(p,\mu,d)}
 {\Sigma\step{\operatorname{start}(p,o)}
  \Sigma[T(p)\mapsto(E[x:=o],\rho,d)]}.
\]
当前命令保持为 $x:=o$，直到操作的结果被写入局部环境。
读取和写入每次访问一个尚未处理的地址：
\[
\frac{T(p)=(E[x:=o],\rho,\Read(U,B))\quad a\in U\quad\mu(a)=v}
 {\Sigma\step{\operatorname{read}(p,a,v)}
  \Sigma[T(p)\mapsto(E[x:=o],\rho,\Read(U\setminus\{a\},B[a\mapsto v]))]}
\]
\[
\frac{T(p)=(E[x:=o],\rho,\Write(U,V))\quad a\in U\quad V(a)=v}
 {\Sigma\step{\operatorname{write}(p,a,v)}
  \Sigma[\mu\mapsto\mu[a\mapsto v]]
        [T(p)\mapsto(E[x:=o],\rho,\Write(U\setminus\{a\},V))]}.
\]
区域内部地址的处理顺序由调度非确定地选择。地址集合为空时，读取状态
转为 $\Done(B)$，写入状态转为 $\Done(\kw{unit})$；计算状态
$\Compute(k,f,\bar v)$ 转为 $\Done(f(\bar v))$。这些步骤产生
$\operatorname{control}(p)$ 标签并保持共享内存。
随后使用规则
\[
\frac{T(p)=(E[x:=o],\rho,\Done(v))}
 {\Sigma\step{\operatorname{complete}(p,x,v)}
  \Sigma[T(p)\mapsto(E[\kw{skip}],\rho[x\mapsto v],\bot)]}.
\]
因此，同一 pipeline 在当前操作完成之前不能执行下一条命令；其他
pipeline 可以继续执行。

\subsection{事件和锁}

对已声明事件 $\varepsilon=(s,p,q,k)$，源 pipeline 设置事件，目标
pipeline 消费事件：
\[
\frac{T(p)=(E[\kw{setflag}(\varepsilon)],\rho,\bot)}
 {\Sigma\step{\operatorname{set}(p,\varepsilon)}
  \Sigma[F\mapsto F[\varepsilon\mapsto1]]
        [T(p)\mapsto(E[\kw{skip}],\rho,\bot)]}
\]
\[
\frac{T(q)=(E[\kw{waitflag}(\varepsilon)],\rho,\bot)\quad F(\varepsilon)=1}
 {\Sigma\step{\operatorname{wait}(q,\varepsilon)}
  \Sigma[F\mapsto F[\varepsilon\mapsto0]]
        [T(q)\mapsto(E[\kw{skip}],\rho,\bot)]}.
\]
重复设置保持标志为 $1$；等待成功时将标志置为 $0$。标志为 $0$ 时等待
没有可用转移。根据顺序语义，执行设置事件时，源 pipeline 的前序操作
已经完成。

令 $p\in\operatorname{Participants}(\ell)$。获取和释放规则为
\begin{gather*}
\frac{T(p)=(E[\kw{lock}(\ell)],\rho,\bot)\quad L(\ell)=\bot}
 {\Sigma\step{\operatorname{acquire}(p,\ell)}
  \Sigma[L\mapsto L[\ell\mapsto p]]
        [T(p)\mapsto(E[\kw{skip}],\rho,\bot)]}\\[1em]
\frac{T(p)=(E[\kw{unlock}(\ell)],\rho,\bot)\quad L(\ell)=p}
 {\Sigma\step{\operatorname{release}(p,\ell)}
  \Sigma[L\mapsto L[\ell\mapsto\bot]]
        [T(p)\mapsto(E[\kw{skip}],\rho,\bot)]}.
\end{gather*}
锁被持有时，其他参与者的获取步骤不可用。事件未声明、事件端点不匹配、
锁未声明、参与权限不足或者释放者不持有锁时进入错误状态。

\subsection{全局调度、错误与终止}

全局步骤非确定地选择一个 $p\in\mathcal P$，执行上述一个可用步骤。
局部控制、操作进度、内存访问及同步操作都属于该 pipeline。
非法命令产生 $\operatorname{error}(p)$ 标签并进入无后继的 $\mathsf{Error}$。
等待标志和等待锁时条件暂未满足属于阻塞。

正常终止定义为所有 pipeline 均满足
$C_p=\kw{skip}$ 且 $d_p=\bot$。未正常终止且没有任何可用步骤的状态是
全局死锁。部分 pipeline 长期等待而其他 pipeline 持续执行的情形由活性
性质描述。

程序语义包含所有合法初态出发的最大有限轨迹和无限轨迹。可选的弱公平条件
要求持续可执行的 pipeline 最终获得调度；基本语义不强制公平。

\section{事件顺序与数据竞争}

对有限标签轨迹 $\tau$，记 $\tau_i$ 为第 $i$ 个事件，$\operatorname{owner}$
返回事件所属 pipeline。基本顺序边包括：
\begin{enumerate}
\item 程序顺序：$i<j$ 且 $\operatorname{owner}(\tau_i)=
  \operatorname{owner}(\tau_j)$。这包括该 pipeline 的读写步骤和同步步骤。
\item 事件同步：$\tau_i=\operatorname{set}(p,\varepsilon)$、
  $\tau_j=\operatorname{wait}(q,\varepsilon)$，$i<j$，中间没有对
  $\varepsilon$ 的设置或等待。
\item 锁同步：$\tau_i=\operatorname{release}(p,\ell)$、
  $\tau_j=\operatorname{acquire}(q,\ell)$，$i<j$，中间没有对 $\ell$
  的获取。
\end{enumerate}
happens-before 是这些边的传递闭包，记为 $\mathit{hb}$。两个不同位置
$i,j$ 构成冲突，当且仅当它们访问同一地址且至少一个是写入。数据竞争为
\[
\operatorname{Race}_\tau(i,j)\iff i\ne j\land
\operatorname{Conflict}_\tau(i,j)\land
\neg(i\mathrel{\mathit{hb}}j)\land\neg(j\mathrel{\mathit{hb}}i).
\]
同一 pipeline 的任意两个先后访问都由程序顺序连接。区域访问可以与其他
pipeline 的访问交错，因而可以产生竞争和混合读取结果。

\section{可选的顺序锁策略}

普通锁语义不指定竞争者的获取次序。可选扩展使用
$\operatorname{order}:\mathsf{Lock}\times\mathbb N\to\mathsf{Pipe}$，
以及每把锁的轮次 $\operatorname{turn}(\ell)$。初始轮次为 $0$；
获取时除了普通规则的前提，还要求
$\operatorname{order}(\ell,\operatorname{turn}(\ell))=p$；
成功释放后将轮次加一。每个轮次的选定 pipeline 必须属于该锁的参与者集合。
忽略轮次状态后，扩展的每个步骤均为普通语义的步骤。无限轮转顺序
$\operatorname{order}(\ell,n)=n\bmod 2$ 可用于参与者为 $0,1$ 的锁。

\section{形式化范围与参考}

Rocq 定义位于 \code{rocq/CoreLanguage/}，包含语法、状态、带标签小步关系、
执行轨迹、happens-before、数据竞争、顺序锁和具体执行推导。
Rocq 证明覆盖操作启动所需的空闲条件、完成时的结果状态、计算类别检查、
正常终止状态没有后继、happens-before 的顺序性质和顺序锁性质。
实际硬件的缓存传播、队列容量、执行时延以及硬件程序翻译到本语言的
refinement 需要针对具体目标建立。将安全结论传递到硬件时还需证明访存、
别名、同步作用域、执行类别和事件顺序得到保持。

事件的二值消费含义参考 Ascend 文档~\cite{ascend-flag,ascend-architecture}；
HIVM 同步操作的定义可见仓库中的
\code{vendor/ascendnpu-ir/upstream/bishengir/include/bishengir/Dialect/HIVM/IR/HIVMSynchronizationOps.td}。
锁的 acquire/release 含义参考 Triton 原子操作接口~\cite{triton-cas,triton-xchg}，
具体翻译由目标平台给出。

\begin{thebibliography}{9}
\bibitem{ascend-flag} Ascend C，\emph{SetFlag/WaitFlag}。
\url{https://asc.gitcode.com/api/SIMD-API/basic_api/sync_control/intra_core_sync/SetFlag-WaitFlag.html}。
\bibitem{ascend-architecture} Ascend C，\emph{Basic Architecture}，CANN 9.0。
\url{https://www.hiascend.com/document/detail/en/CANNCommunityEdition/900/programug/Ascendcopdevg/atlas_ascendc_10_0008.html}。
\bibitem{triton-cas} Triton，\emph{triton.language.atomic\_cas}。
\url{https://triton-lang.org/main/python-api/generated/triton.language.atomic_cas.html}。
\bibitem{triton-xchg} Triton，\emph{triton.language.atomic\_xchg}。
\url{https://triton-lang.org/main/python-api/generated/triton.language.atomic_xchg.html}。
\end{thebibliography}

\clearpage
\section{具体执行例子}

\subsection{同一 pipeline 的操作顺序}

\begin{lstlisting}
parallel {
  pv:
    x := computation[Vector, add](2, 3);
    y := computation[Vector, add](x, 10);
}
\end{lstlisting}
第一条计算执行期间，$d_{p_v}$ 依次为
$\Compute(\mathsf{Vector},\mathrm{add},[2,3])$ 和 $\Done(5)$。
结果写入 $x$ 后，第一条命令变成 $\kw{skip}$；消除顺序组合中的
$\kw{skip}$ 后，第二条计算才能求值 $x$ 并得到 $y=15$。

\subsection{Cube 和 Vector pipeline 的并行计算}

设 $\kind(p_c)=\mathsf{Cube}$、$\kind(p_v)=\mathsf{Vector}$，
$A=\bigl(\begin{smallmatrix}1&2\\3&4\end{smallmatrix}\bigr)$，
$B=\bigl(\begin{smallmatrix}5&6\\7&8\end{smallmatrix}\bigr)$。
\begin{lstlisting}
parallel {
  pv: z := computation[Vector, vadd]([1, 2], [10, 20]);
  pc: c := computation[Cube, matmul](A, B);
}
\end{lstlisting}
一条合法执行依次启动 Vector 计算、启动 Cube 计算、完成 Vector 计算并
写入 $z=[11,22]$，再完成 Cube 计算并写入
$c=\bigl(\begin{smallmatrix}19&22\\43&50\end{smallmatrix}\bigr)$。
两次计算分别保存在所属 pipeline 的当前操作状态中。
将 Cube 计算放在 $p_v$ 中会因执行类别不匹配进入错误状态。

\subsection{区域读取与写入交错}

初始内存满足 $\mu(0)=\mu(1)=0$，两个 pipeline 均具有相应访问权限：
\begin{lstlisting}
parallel {
  p0: result := read({0, 1});
  p1: unit := write({0, 1}, {0: 41, 1: 42});
}
\end{lstlisting}
一条合法执行先由 $p_0$ 读取地址 $0$ 的值 $0$，再由 $p_1$ 写入地址
$0,1$，最后由 $p_0$ 读取地址 $1$ 的值 $42$。返回的映射为
$\{0\mapsto0,1\mapsto42\}$，最终内存为 $\{0\mapsto41,1\mapsto42\}$。
这些逐地址步骤对应 \code{MemoryExamples.v} 中经过证明的执行轨迹。

\subsection{事件与临界区}

设 $\varepsilon$ 的源为 $p_0$、目标为 $p_1$，初始标志为 $0$：
\begin{lstlisting}
parallel {
  p0: unit := write({a}, {a: 41}); setflag(epsilon);
  p1: waitflag(epsilon); result := read({a});
}
\end{lstlisting}
$p_0$ 完成写入后才能设置事件；$p_1$ 消费事件后读取地址 $a$。
当没有其他写入时，读取结果为 $41$。

两个 pipeline 也可以使用同一把锁保护地址 $a$：
\begin{lstlisting}
parallel {
  p0: lock(l); unit := write({a}, {a: 5}); unlock(l);
  p1: lock(l); result := read({a}); unlock(l);
}
\end{lstlisting}
每个 pipeline 的访存操作完成后才执行释放锁。锁的获取顺序决定两个
临界区的先后顺序。

\subsection{循环等待}

设事件 $\varepsilon_{01}$ 从 $p_0$ 发往 $p_1$，$\varepsilon_{10}$
从 $p_1$ 发往 $p_0$，初始标志均为 $0$：
\begin{lstlisting}
parallel {
  p0: waitflag(e10); setflag(e01);
  p1: waitflag(e01); setflag(e10);
}
\end{lstlisting}
两个 pipeline 的当前等待都无法执行，程序未终止且不存在其他可用步骤，
因此初始状态满足全局死锁定义。

\end{document}
